\lecture{5.6}{Super Learner Ensembles}
\label{lec:5-6}

\subsection{Weighted Combinations of Estimators}

The discrete Super Learner chooses one candidate algorithm. But different algorithms
may capture different features of the data. A regression tree may capture threshold
effects, a generalized additive model may capture smooth nonlinearities, and a lasso
model may be stable in moderately high dimensions. Instead of choosing exactly one,
we can combine them.

Let $\Qbar_{n,1},\ldots,\Qbar_{n,K}$ be fitted candidate outcome regressions. A
weighted ensemble has the form
\[
\Qbar_{n,\omega}(a,w)
=
\sum_{k=1}^K \omega_k \Qbar_{n,k}(a,w),
\]
where $\omega=(\omega_1,\ldots,\omega_K)$ is a vector of weights.

\subsection{Convex Combinations}

Super Learner commonly uses convex combinations:
\[
\omega_k\ge 0,\qquad \sum_{k=1}^K\omega_k=1.
\]
The set of all such weights is the simplex
\[
\Delta_K=\left\{\omega\in\R^K:\omega_k\ge 0,\ \sum_{k=1}^K\omega_k=1\right\}.
\]

Convex combinations have a simple interpretation. If $\omega_3=0.70$ and
$\omega_5=0.30$, the ensemble prediction is 70\% from algorithm 3 and 30\% from
algorithm 5. Algorithms with weight zero do not contribute to the final prediction.

\begin{figure}[H]
\centering
\begin{tikzpicture}[
  node distance=1.15cm,
  box/.style={draw, rounded corners=2pt, minimum width=2.7cm, minimum height=.75cm},
  arr/.style={-{Latex[length=2.5mm]}, thick}
]
\node[box] (a1) {Algorithm 1\\ $\Qbar_{n,1}$};
\node[box, below=of a1] (a2) {Algorithm 2\\ $\Qbar_{n,2}$};
\node[box, below=of a2] (a3) {Algorithm 3\\ $\Qbar_{n,3}$};
\node[box, right=2.2cm of a2] (ens) {Ensemble\\ $\sum_k\omega_k\Qbar_{n,k}$};
\draw[arr] (a1) -- node[above, sloped] {$\omega_1$} (ens);
\draw[arr] (a2) -- node[above] {$\omega_2$} (ens);
\draw[arr] (a3) -- node[below, sloped] {$\omega_3$} (ens);
\end{tikzpicture}
\caption{A Super Learner ensemble combines candidate predictions using learned weights.}
\label{fig:ch5-ensemble}
\end{figure}

\subsection{Choosing Weights by Cross-Validated Risk Minimization}

The Super Learner chooses $\omega$ to minimize cross-validated risk. For each fold
$v$, each candidate algorithm is trained on the training folds and predicts the
validation fold. For a proposed weight vector $\omega$, the ensemble prediction in
the validation fold is
\[
\Qbar_{n,\omega,v}(a,w)
=
\sum_{k=1}^K\omega_k\Qbar_{n,k,v}(a,w).
\]
The cross-validated risk of the ensemble is
\[
R_{CV,n}(\omega)
=
\frac{1}{V}\sum_{v=1}^V
\frac{1}{|\mathcal V_v|}
\sum_{i\in\mathcal V_v}
L(\Qbar_{n,\omega,v})(W_i,A_i,Y_i).
\]
The selected weights are
\[
\omega_{CV}
=
\argmin_{\omega\in\Delta_K}R_{CV,n}(\omega).
\]
After the weights are selected, each candidate algorithm is refit on the full dataset,
and the final Super Learner is
\[
\Qbar_{n,SL}(a,w)
=
\sum_{k=1}^K\omega_{CV,k}\Qbar_{n,k}^{\mathrm{full}}(a,w).
\]

\subsection{Oracle Inequalities}

The theoretical appeal of Super Learner comes from oracle inequalities. An oracle is
an imaginary procedure that knows the true risks of all candidates and chooses the
best one. The analyst cannot implement the oracle because true risk depends on the
unknown data-generating distribution.

An oracle inequality states, under regularity conditions, that the cross-validation
selector or ensemble performs nearly as well as the oracle, up to a term that becomes
small as sample size grows. A simplified version of the message is
\[
\text{risk of Super Learner}
\le
\text{risk of oracle choice}
+
\text{small remainder}.
\]
The remainder depends on sample size and on the size of the candidate library. This
supports the practical strategy of including a diverse pre-specified library rather than
trying to guess a single correct model.

\begin{keyidea}
The oracle is the best procedure in the chosen library for the data-generating
distribution. Super Learner uses cross-validation to approximate that oracle without
knowing the truth.
\end{keyidea}

\subsection{Model Stacking}

Combining candidate predictions is often called stacking. The idea is to treat the
out-of-fold predictions from the candidate algorithms as inputs to a second-stage
learner, sometimes called the metalearner. In the convex Super Learner, the
metalearner is a constrained regression that finds nonnegative weights summing to
one.

Stacking differs from ordinary model averaging based on subjective preference. The
weights are chosen to minimize a pre-specified cross-validated loss. This makes the
ensemble a fully specified estimator rather than a post-hoc compromise among models.

\subsection{Super Learner vs Single Best Algorithm}

The discrete Super Learner is guaranteed to choose one candidate. The ensemble Super
Learner can choose one candidate by putting all weight on it, but it can also combine
multiple candidates when that lowers cross-validated risk. Thus the ensemble contains
the discrete choice as a special case.

\begin{estimationbox}
Ensemble Super Learner workflow:
\begin{enumerate}
\item Pre-specify candidate algorithms and a valid loss function.
\item Generate out-of-fold predictions for every algorithm.
\item Choose convex weights minimizing cross-validated risk.
\item Refit all candidate algorithms on the full dataset.
\item Combine full-data predictions using the selected weights.
\item Use the resulting nuisance function estimate inside the causal estimator.
\end{enumerate}
\end{estimationbox}
